% META: {"id": "1NSI-APT-EVAL-A-corrige", "chapitre": "1NSI-ALGO-PARCOURS-TRIS", "type_objet": "corrige_evaluation", "evaluation_ref": "1NSI-APT-EVAL-A", "capacites": ["P-ALGO-01A", "P-ALGO-01B", "P-ALGO-02A", "P-ALGO-02B", "P-ALGO-02C", "P-ALGO-02D"], "mode_creation": "ex_nihilo", "sources_inspiration": [], "status": "needs_review"}
% BEGIN-VERIFY
% def recherche_occurrence(tableau, valeur):
%     for i in range(len(tableau)):
%         if tableau[i] == valeur:
%             return i
%     return -1
% def maximum(tableau):
%     maxi = tableau[0]
%     for x in tableau[1:]:
%         if x > maxi:
%             maxi = x
%     return maxi
% def moyenne(tableau):
%     return sum(tableau) / len(tableau)
% t = [22, 15, 8, 30, 11]
% assert recherche_occurrence(t, 30) == 3
% assert maximum(t) == 30
% assert moyenne(t) == 17.2
% def tri_insertion(tableau):
%     tableau = tableau[:]
%     for i in range(1, len(tableau)):
%         valeur = tableau[i]
%         j = i - 1
%         while j >= 0 and tableau[j] > valeur:
%             tableau[j + 1] = tableau[j]
%             j -= 1
%         tableau[j + 1] = valeur
%     return tableau
% assert tri_insertion([9, 2, 6, 4]) == [2, 4, 6, 9]
% END-VERIFY
\begin{corrige}{1NSI-APT-EVAL-A-EX1}
\textbf{1.} La valeur $30$ se trouve à l'indice $3$.

\textbf{2.} Le maximum est $30$. L'accumulateur ne peut pas être initialisé à $0$ : sur
un tableau dont toutes les valeurs seraient négatives, aucune ne dépasserait $0$, et la
fonction renverrait $0$ — une valeur qui n'appartient pas au tableau. L'initialisation doit
donc venir du tableau lui-même, avec \lstinline{plus_grand = tableau[0]}.

\textbf{3.} $\dfrac{22+15+8+30+11}{5}=\dfrac{86}{5}=17{,}2$.
\end{corrige}

\begin{corrige}{1NSI-APT-EVAL-A-EX2}
\textbf{1.} Tableau initial : \lstinline{[9, 2, 6, 4]}.
\begin{itemize}
  \item Après $i=1$ (insertion de $2$) : \lstinline{[2, 9, 6, 4]}
  \item Après $i=2$ (insertion de $6$) : \lstinline{[2, 6, 9, 4]}
  \item Après $i=3$ (insertion de $4$) : \lstinline{[2, 4, 6, 9]}
\end{itemize}

\textbf{2.} Invariant : \emph{au début de chaque itération $i$, le sous-tableau
\lstinline{tableau[0:i]} contient les $i$ premiers éléments d'origine, rangés en ordre
croissant.} Borner la zone est indispensable : pendant l'exécution le tableau entier n'est
pas trié. Vérification — après $i=1$ : \lstinline{[2, 9]} trié ; après $i=2$ :
\lstinline{[2, 6, 9]} trié ; après $i=3$ : \lstinline{[2, 4, 6, 9]} trié.

\textbf{3.} Le coût dans le pire cas est \textbf{quadratique}, $O(n^2)$.
\end{corrige}

\begin{corrige}{1NSI-APT-EVAL-A-EX3}

\textbf{1.} Les deux lignes manquantes :
\begin{python}
def tri_selection(tableau):
    tableau = tableau[:]
    n = len(tableau)
    for i in range(n - 1):
        indice_min = i                    # (a)
        for j in range(i + 1, n):
            if tableau[j] < tableau[indice_min]:
                indice_min = j            # (b)
        tableau[i], tableau[indice_min] = tableau[indice_min], tableau[i]
    return tableau
\end{python}
La ligne (a) doit rester \textbf{avant} la boucle sur $j$ : placée à l'intérieur, elle
réinitialiserait le candidat à chaque tour et l'algorithme ne retiendrait plus le minimum
global de la zone.

\textbf{2.} L'invariant du tri par sélection, au début de l'itération $i$ :
\begin{itemize}
  \item (a) \lstinline{tableau[0:i]} est trié ;
  \item (b) chacun de ses éléments est inférieur ou égal à chacun des éléments de
    \lstinline{tableau[i:n]}.
\end{itemize}

\textbf{3.} C'est la clause \textbf{(b)} qui est fausse pour le tri par insertion. Sur le
tableau de l'exercice 2, au début du tour $i=3$, la zone triée est \lstinline{[2, 6, 9]} et
il reste \lstinline{[4]} : on aurait besoin de $9 \leqslant 4$, ce qui est faux. Le $4$ devra
d'ailleurs traverser une partie de la zone déjà triée — exactement ce que la clause (b)
interdit dans le tri par sélection, où un élément placé est définitif.

\end{corrige}
